\documentclass[12pt]{article}
\usepackage[utf8]{inputenc}
\usepackage[T1]{fontenc}
\usepackage[margin=1in]{geometry}
\usepackage{setspace}
\usepackage{xcolor}
\usepackage[colorlinks=true, urlcolor=blue!60!black, linkcolor=blue!60!black]{hyperref}
\setlength{\parindent}{0em}
\setlength{\parskip}{0.45em}
\renewcommand{\baselinestretch}{1.0}
\begin{document}
\begin{center}
{\Large \textit{Fletcher v. Peck}, 10 U.S. 87 (1810)\par}
\vspace{0.3em}
{\normalsize Case Brief --- One-Page Edition\par}
\end{center}

\textbf{Memory Jogger.} The Yazoo land fraud case: a state land grant is a ``contract'' Georgia could not rescind, the first state statute the Court struck down (\href{https://www.courtlistener.com/opinion/84935/fletcher-v-peck/}{10 U.S. 87}).

\textbf{Facts.} A bribed Georgia legislature sold 35 million Yazoo acres in 1795; its successor voided the grant in 1796 after tracts had passed to innocent purchasers. Peck sold land to Fletcher with title covenants, and Fletcher sued for breach in a ``feigned'' test case.

\textbf{Procedural History.} The federal circuit court in Massachusetts ruled for Peck; Fletcher brought a writ of error.

\textbf{Votes.} Affirmed, 4--1 on the reasoning: Marshall for the Court; Johnson concurred in part and dissented in part; two Justices did not participate.

\textbf{Holding.} A completed state land grant is an executed contract under Article~I, \S~10; once title vests in bona fide purchasers, a later legislature cannot rescind it, and courts will not void statutes by probing legislators' motives.

\textbf{Reasoning.} Courts cannot count tainted votes, and equity protects purchasers without notice. A grant extinguishes the grantor's right, and the Clause draws no line between private and public contracts. Johnson relied instead on natural law, ``a principle which will impose laws even on the deity,'' and dissented on Indian title.

\textbf{Examples.} \emph{Same-Side:} revoking a university charter (\textit{Dartmouth College}). \emph{Opposite-Side:} prospectively repealing an unvested subsidy. \emph{Fence-Sitter:} suing in equity to cancel a bribed grant still held by the wrongdoers.

\textbf{Critique.} Textualists say the Clause targeted private debtor relief; progressives see a corrupt giveaway shielded from democratic reversal.

\textbf{Key Quotation.} ``A contract executed, as well as one which is executory, contains obligations binding on the parties. A grant, in its own nature, amounts to an extinguishment of the right of the grantor, and implies a contract not to reassert that right. A party is, therefore, always estopped by his own grant.'' --- Chief Justice Marshall, Opinion of the Court, 10 U.S. at 137.

\end{document}
